% Осень26!A28; нумерация внутри столбца без пустых строк.
% Источник: лекция 5, §2 (строка 68)
% Источник: лекция 5, §3 (строка 115)
\ticket{25}{Теорема об измеримости пределов и супремумов}
\label{ticket:25}

\begin{theorem}
Пусть $f_n:X\to\overline{\mathbb R}$ --- измеримые функции.
Тогда измеримы
\[
  \sup_n f_n,\quad\inf_n f_n,\quad
  \limsup_{n\to\infty}f_n=\inf_n\sup_{k\ge n}f_k,\quad
  \liminf_{n\to\infty}f_n=\sup_n\inf_{k\ge n}f_k.
\]
Если $f_n(x)\to f(x)$ в каждой точке $X$, то $f$ измерима.
\end{theorem}

\begin{proof}
Измеримость можно проверять по любому из четырёх типов множеств Лебега, поскольку
\[
  \{u>a\}=X\setminus\bigcap_{j=1}^{\infty}\{u<a+1/j\},
\]
\[
  \{u<a\}=\bigcup_{j=1}^{\infty}\bigl(X\setminus\{u>a-1/j\}\bigr).
\]
Для супремума и инфимума имеем
\[
  \{\sup_n f_n>a\}=\bigcup_n\{f_n>a\},
  \qquad
  \{\inf_n f_n<a\}=\bigcup_n\{f_n<a\}.
\]
Правые части измеримы как счётные объединения измеримых множеств.
Функции $\sup_{k\ge n}f_k$ и $\inf_{k\ge n}f_k$ измеримы по доказанному.
Ещё одно взятие инфимума или супремума даёт измеримость верхнего и нижнего пределов.
При существовании поточечного предела он совпадает с верхним пределом и потому измерим.
\end{proof}
